% Discussion: the result-dependent paragraphs are written from paper/data once the campaigns are complete.
% generated / written once the campaigns are complete


\textbf{Limitations.} The models are small (12k--260k parameters) and the conclusions about architectures are drawn at
this scale; Narang et al.\ \cite{narang2021transformer} show that architectural conclusions do not always transfer across
scales, and outlier behaviour changes with scale \cite{dettmers2022llmint8}. On TinyStories every model is trained once
(one seed, one learning rate) for compute reasons. All weights and the KV cache are on chip, which bounds the model size on
the XC7Z020 at roughly the size of stories260K; larger models need weights streamed from DDR, which changes the
performance model from compute-bound to bandwidth-bound. The attention unit assumes a head size of 8. Resource numbers
for every model come from open-source synthesis (Yosys) without timing; post-route utilization, timing and power for the
one-bitstream engine come from Vivado (\S\ref{sec:res-board}). Energy is measured at the board input with a USB meter,
which includes the PS and the DDR; we therefore report dynamic power above idle.

\textbf{Reproducibility.} The repository contains the training campaigns (seeds, learning-rate sweeps, logs), the
compiler, the integer reference, the generated memories and golden vectors of every model, the RTL and testbenches, the
Vivado scripts for the overlay and the out-of-context sweep, the PYNQ driver and measurement scripts, and the scripts that
generate every table and figure of this paper from the raw results. A continuous-integration job re-runs the bit-exact
regression on every commit.
